\documentclass[conference]{IEEEtran}
\usepackage{flushend}
\usepackage{xurl}
\setlength{\emergencystretch}{2em}
\title{A Causal Infrastructure Advisory for Freeway Queue Encounters}
\author{Arjun Vijay Prakash \\ \small IEEE Intelligent Vehicles Symposium 2027 | Workshop
working draft}

\begin{document}

\maketitle
\begin{abstract}
An infrastructure warning can give a moving vehicle more time to react to a queue, but stale or
weakly localized warnings can reach vehicles that will never encounter it. CREST sends a
direction, forecast horizon, expiry, and queue-tail location interval. A receiver can check
availability and the projected position of its probe against the warned intervals. In a
retrospective replay of nine five-minute blocks from one US-101 recording, the gated method
reduced false warnings from 9,404 to 136 across 9,639 vehicle-and-forecast-origin pairs at a
7.5-second delay. It detected only 9 of 235 proxy encounters. At 30 seconds the advisory expired
before delivery on every pair. The replay identifies a recall failure that limits the
vehicle-facing utility of this advisory.
\end{abstract}

\section*{Advisory and causal replay}
The queue head uses directed history and only readings available by the forecast origin. It
assigns queue probabilities to cells. The advisory publishes contiguous confident intervals,
with a separate interval for the upstream queue tail. The message carries issuance and expiry
times and an unavailable state. The consumer projects the vehicle's position from its origin
position and speed. Later positions are reserved for encounter scoring, so a probe's realized
motion cannot determine whether it receives a warning. The deployment concept is read-only and
the experiment uses replay.

\section*{Development setup}
The model trained on nine five-minute blocks from two I-80 recordings and tested on nine
nonoverlapping blocks from one US-101 recording. It forecast two 15-second steps ahead.
We crossed 15\%, 35\%, and 70\% probe penetration with 0-, 7.5-, and 30-second delivery delays.
The reference uses contiguous low-speed cells calculated from NGSIM trajectories.

We compared an infrastructure advisory that requires a fresh, localized interval intersecting
the projected path, an ungated comparator that thresholds probabilities along that path, and a
probe-speed rule. A proxy encounter starts outside a queue and ends inside one at the forecast
horizon. Warnings on pairs without a proxy encounter count as false warnings.

\section*{Development results}
At 35\% penetration and 7.5 seconds of delay, the gated and ungated infrastructure methods
detected 9 and 235 encounters and produced 136 and 9,404 false warnings, respectively.
The probe-speed rule detected 55 encounters, gave 3,172 false warnings, and withheld output
on 6,335 pairs. Median lead time was 22.5 seconds among each method's detections because
the replay fixes the forecast horizon. At 30 seconds the gated advisory withheld all
9,639 pairs. The blocks share one test recording, so the counts describe this transfer and
cannot establish performance across independent recordings.

\section*{References}
[1] Federal Highway Administration. Next Generation Simulation vehicle trajectories and
supporting data.
\url{https://catalog.data.gov/dataset/next-generation-simulation-ngsim-vehicle-trajectories-and-supporting-data}

[2] CREST measured development receipt.
\url{https://github.com/ArjunCodess/CREST/blob/main/research_manuscript/generated/roadmap_development.json}
\end{document}
